% Fragmento integrable. Preambulo anfitrion: amsmath, amssymb, longtable, hyperref.
\section{Transverse entry and local scaling}
\label{hmtfg:fragmento:entrada-transversal-y-escalado-local}

\subsection{Provenance and result}
\label{hmtfg:escalado:1}

APP evaluates addition and multiplication on the two sheets of the lattice, preserving quotient and remainder. TRIT determines regime and orientation. TPK transport selects, transports and updates the state with its carry, boundary and memory records. Its readers act on the same discrete structure of the continuum. The prolongation \texttt{w6→w12→w18→w24→w30→R36→G9} preserves this provenance and distinguishes phase return from memory increment.

The critical reader of the dodecaphase wheel, constructed in the operator antecedents and transported by compatible refinement, produces
\[
\phi_j=\frac{(3j-1)\pi_{\mathrm{HMT}}}{18},\qquad
w_j=\frac1{12},\qquad
k_j=\sqrt{\frac2{\sum_{\ell=1}^{12}w_\ell\phi_\ell^2}}\,\phi_j
=\frac{2(3j-1)}{\sqrt{899}}.
\]
The coordinate \(\pi_{\mathrm{HMT}}\) retains its APP–TRIT–TPK origin; its common factor cancels in this second moment. The resulting reader is
\begin{equation}
u_0(x)=\frac1{12}\sum_{j=1}^{12}(1-\cos k_jx),\qquad
f_\lambda(x)=1-\lambda u_0(x).
\label{hmtfg:escalado:eq:1}
\end{equation}
The focal sector is the uniform one, \(\eta=0\). The preceding extension separately retains the weighted family \(|\eta|\le1/40\). The new numerical bounds correspond to \eqref{hmtfg:escalado:eq:1}.

The family is even and entire, with \(f_\lambda(0)=1\), a quadratic critical point and second moment \(\sum_jw_jk_j^2=2\). The parameter \(\lambda\) ranges over a family downstream of the HMT generator. The target Feigenbaum constant remains outside the inputs of \eqref{hmtfg:escalado:eq:1}.

\textbf{Main result.} There exists a unique
\begin{equation}
\lambda_*\in[1.874038,1.874039]
\label{hmtfg:escalado:eq:2}
\end{equation}
for which the fourth normalized return of \eqref{hmtfg:escalado:eq:1} belongs to the local stable leaf of the real doubling fixed point. The crossing of the family with that leaf is transverse, with quantified separation. For every \(n\ge0\),
\begin{equation}
\boxed{\|R^{4+n}f_{\lambda_*}-g\|_{\mathcal A}
<0.001666\,(0.83996)^n.}
\label{hmtfg:escalado:eq:3}
\end{equation}
The local superstable cascade associated with this crossing has parameters \(\mu_j\) and constants \(C\ne0\), \(0<\theta<1\), such that
\begin{equation}
\mu_j-\lambda_*=C\delta_{\mathrm F}^{-j}\bigl[1+O(\theta^j)\bigr],
\qquad
\frac{\mu_{j-1}-\mu_{j-2}}{\mu_j-\mu_{j-1}}\longrightarrow\delta_{\mathrm F}.
\label{hmtfg:escalado:eq:4}
\end{equation}
Here \(\delta_{\mathrm F}>1\) is the unstable eigenvalue of the doubling operator. Selection of \(\mu_j\) is performed through the local returns of section~\ref{hmtfg:escalado:8}. Identification with the global selector is proved in sections~\ref{hmtfg:escalado:15}--\ref{hmtfg:escalado:17}: all global terms retain the combinatorics and the centers agree at every sufficiently high level when their periods are matched.


\subsection{Analytic chart and effective return}
\label{hmtfg:escalado:2}

The chart of even functions is used
\[
s=x^2,\qquad z=\frac{s-1}{5/2},\qquad
f(x)=1-x^2V(z),\qquad V(z)=\sum_{j\ge0}v_jz^j.
\]
The coordinates of the Banach space are
\begin{equation}
u=10v_0,\qquad \nu=(v_1,v_2,\ldots),\qquad
\|(\Delta u,\Delta\nu)\|_{\mathcal A}
=|\Delta u|+\sum_{j\ge1}|\Delta v_j|.
\label{hmtfg:escalado:eq:5}
\end{equation}
In particular, the constant coefficient of \(V\) has weight ten. This normalization agrees with the Hertling–Spandl chart \(u/10\), section 2 of the source.

Let \(a=-f(1)=v_0-1\). For admissible real returns,
\begin{equation}
(Rf)(x)=-\frac{f(f(-ax))}{a}.
\label{hmtfg:escalado:eq:6}
\end{equation}
The program evaluates an exact factorization of \eqref{hmtfg:escalado:eq:6}. Defining
\[
S=\frac{a^2s-1}{5/2},\qquad b=V(S),\qquad
h=1-a^2sb,\qquad t=\frac{h^2-1}{5/2},
\]
\[
\operatorname{shift}V(z)=\sum_{j\ge0}v_{j+1}z^j,
\]
one obtains
\begin{equation}
\boxed{V_R=a\,b\,(h+1)\left[V(t)+\frac25\operatorname{shift}V(t)\right].}
\label{hmtfg:escalado:eq:7}
\end{equation}
To verify it, one uses \(v_0=1+a\),
\(h^2-1=-a^2sb(h+1)\) and
\(V(t)-v_0=t\operatorname{shift}V(t)\). Thus,
\(a+1-h^2V(t)=a^2sb(h+1)[V(t)+(2/5)\operatorname{shift}V(t)]\),
which is precisely \(asV_R\).

The initial moments are rational:
\begin{equation}
\frac1{12}\sum_jk_j^{2m}
=\frac1{12}\left(\frac4{899}\right)^m
\sum_{j=1}^{12}(3j-1)^{2m}.
\label{hmtfg:escalado:eq:8}
\end{equation}
Therefore both \eqref{hmtfg:escalado:eq:7} and the initial construction use rational intervals; the initial trigonometric functions are represented by their series with an outer remainder bound.


\subsection{External bounds used as subsequent recognition}
\label{hmtfg:escalado:3}

Once \eqref{hmtfg:escalado:eq:1} has been constructed, Lanford’s hyperbolicity theorem is used as an analytic certificate for the operator \eqref{hmtfg:escalado:eq:6}. Let \(g_0\) be the published central polynomial and \(g\) the fixed point. The estimates used are
\[
\|g-g_0\|_{\mathcal A}<\varepsilon_0=4\,10^{-5},
\]
and, in the ball \(\|f-g_0\|_{\mathcal A}<0.01\), the blocks of \(DR\), written \(R=(U,V)\), satisfy
\begin{equation}
U_u\ge a_0=4.521,\quad
\|U_\nu\|\le b_0=0.560,\quad
\|V_u\|\le c_0=0.756,\quad
\|V_\nu\|\le d_0=0.719.
\label{hmtfg:escalado:eq:9}
\end{equation}
The fixed point has a unique unstable direction, with positive real eigenvalue \(\delta_{\mathrm F}>1\); the rest of the spectrum lies inside the unit disk. The figures in \eqref{hmtfg:escalado:eq:9} come from the published estimates, not from a new reproduction of Lanford’s proof in this dossier. The coefficients of \(g_0\) are explicitly incorporated into the program’s central reference polynomial, in accordance with table 2 of Hertling–Spandl.

Sources: \href{https://repo-archives.ihes.fr/A_GARDER/20180409/Test/I_Prepublications/LANFORD/1968-2013/P_81_17/P_81_17_web.pdf}{Lanford, hyperbolicity estimates}; \href{https://arxiv.org/pdf/1410.3277}{Hertling–Spandl, analytic chart and table 2}.


\subsection{Rational entry and direction certificate}
\label{hmtfg:escalado:4}

The renormalization entry certificate divides \eqref{hmtfg:escalado:eq:2} into sixteen equal rational intervals. It propagates four returns of \eqref{hmtfg:escalado:eq:7} and the exact derivative with respect to \(\lambda\), using differentiation of the composition and outward rounding to one hundred decimal places.

Each series consists of a polynomial of degree 32 with interval coefficients and an arbitrary analytic error \(E\), \(\|E\|_1\le\epsilon\). This error may also modify low-order coefficients. For \(\|q\|_1<1\),
\begin{equation}
\|E\circ q\|_1\le\epsilon,\qquad
\|E'\circ q\|_1\le\frac{\epsilon}{(1-\|q\|_1)^2},\qquad
\|\operatorname{shift}E\|_1\le\epsilon.
\label{hmtfg:escalado:eq:10}
\end{equation}
The products include polynomial–error and error–error terms. The initial tail is bounded by its first absolute term and a geometric ratio, using \(\|1+(5/2)z\|_1=7/2\). Each composed argument remains strictly inside the unit ball of this algebra.

For \(\Gamma(\lambda)=R^4f_\lambda=(u_\lambda,\nu_\lambda)\), the certificate gives the following conservative bounds, uniform in \eqref{hmtfg:escalado:eq:2}:

\begingroup\small
\begin{longtable}{p{0.465\linewidth}p{0.465\linewidth}}
\hline
Quantity & Certified outer bound \\
\hline
\(\|\Gamma(\lambda)-g_0\|_{\mathcal A}\) & \(<0.004993128\) \\
\(\|\nu_\lambda-\nu_{g_0}\|_1\) & \(<0.001396077\) \\
\(u'_\lambda\) & \(>3821.289115\) \\
\(\|\nu'_\lambda\|_1\) & \(<551.757000\) \\
\(\|\nu'_\lambda\|_1/u'_\lambda\) & \(<0.144386<1/4\) \\
Analytic error of the function & \(<8.482\,10^{-19}\) \\
Analytic error of the tangent & \(<4.569\,10^{-15}\) \\
\hline\end{longtable}
\endgroup

The complete decimals, the sixteen subintervals and every intermediate return are retained in the rational entry certificate.

The cone \(\|\dot\nu\|_1\le|\dot u|/4\) is invariant under \eqref{hmtfg:escalado:eq:9}:
\[
|\dot u_{\rm nuevo}|\ge(4.521-0.560/4)|\dot u|
=4.381|\dot u|,
\]
\begin{equation}
\|\dot\nu_{\rm nueva}\|_1\le(0.756+0.719/4)|\dot u|
=0.93575|\dot u|<\tfrac14(4.381)|\dot u|.
\label{hmtfg:escalado:eq:11}
\end{equation}
The certified parameter direction therefore enters a uniformly expanding cone of the operator.


\subsection{Quantitative construction of the stable leaf}
\label{hmtfg:escalado:5}

Translate \(g\) to the origin and write the coordinates as \((x,y)=(u-u_g,\nu-\nu_g)\). Fix
\begin{equation}
L=0.16,\qquad r=0.004,\qquad
A=a_0-Lc_0=4.40004,\qquad q=d_0+c_0L=0.83996.
\label{hmtfg:escalado:eq:12}
\end{equation}
The cylinder \(|x|\le Lr\), \(\|y\|_1\le r\) lies inside the ball of \eqref{hmtfg:escalado:eq:9}, since
\((1+L)r+\varepsilon_0=0.00468<0.01\).

Consider the complete space of functions \(\sigma:B_r\to\mathbb R\) with \(\sigma(0)=0\) and Lipschitz constant at most \(L\), equipped with the uniform distance. For each \(y\), define \((\mathcal G\sigma)(y)\) as the root in \([-L\|y\|_1,L\|y\|_1]\) of
\begin{equation}
F_\sigma(x,y)=U(x,y)-\sigma(V(x,y)).
\label{hmtfg:escalado:eq:13}
\end{equation}
On that interval, \(\|V(x,y)\|_1\le q\|y\|_1<r\). As \(x\) increases, \eqref{hmtfg:escalado:eq:9} shows that \(F_\sigma\) grows with slope at least \(A\). At its endpoints,
\[
F_\sigma(L\|y\|_1,y)
\ge[AL-b_0-Ld_0]\|y\|_1,
\]
and the opposite inequality holds at the negative endpoint. The margin is
\begin{equation}
AL-b_0-Ld_0=0.0289664>0.
\label{hmtfg:escalado:eq:14}
\end{equation}
A unique root is obtained. Comparing \eqref{hmtfg:escalado:eq:13} for two values of \(y\),
\[
\operatorname{Lip}(\mathcal G\sigma)
\le\frac{b_0+Ld_0}{A}
=\frac{0.67504}{4.40004}<0.16.
\]
Comparing two graphs,
\begin{equation}
\|\mathcal G\sigma-\mathcal G\widetilde\sigma\|_\infty
\le A^{-1}\|\sigma-\widetilde\sigma\|_\infty,
\qquad A^{-1}<0.228.
\label{hmtfg:escalado:eq:15}
\end{equation}
The contraction theorem produces a unique fixed graph \(x=\sigma_*(y)\). On it,
\begin{equation}
\|y_n\|_1\le q^n\|y_0\|_1,
\qquad
\|(x_n,y_n)\|_{\mathcal A}\le(1+L)q^n\|y_0\|_1.
\label{hmtfg:escalado:eq:16}
\end{equation}
Moreover, the vertical distance \(d(x,y)=x-\sigma_*(y)\) satisfies
\[
|d(R(x,y))|\ge A|d(x,y)|
\]
as long as the orbit remains in the cylinder. Every orbit that remains there indefinitely belongs to the graph. This graph identifies exactly the local stable leaf and provides estimate \eqref{hmtfg:escalado:eq:16}. Its local regularity agrees with the analytic stable manifold of the hyperbolic fixed point.


\subsection{Existence, uniqueness and transversality of the parameter}
\label{hmtfg:escalado:6}

For \(\lambda\) in \eqref{hmtfg:escalado:eq:2}, define
\[
H(\lambda)=u_\lambda-u_g-
\sigma_*(\nu_\lambda-\nu_g).
\]
The certificate ensures \(\|\nu_\lambda-\nu_g\|_1<0.001436077<r\). For \(\lambda_2>\lambda_1\),
\begin{equation}
H(\lambda_2)-H(\lambda_1)
\ge\int_{\lambda_1}^{\lambda_2}
\bigl[u'_t-L\|\nu'_t\|_1\bigr]dt
>3733.007995\,(\lambda_2-\lambda_1).
\label{hmtfg:escalado:eq:17}
\end{equation}
To determine the endpoint signs, \(\|g-g_0\|_{\mathcal A}<\varepsilon_0\) and \(|\sigma_*(y)|\le L\|y\|_1\) are used. Rational point evaluation produces
\begin{equation}
H(1.874038)<-0.0011856110945,
\qquad H(1.874039)>0.0021866053399.
\label{hmtfg:escalado:eq:18}
\end{equation}
By continuity, a root exists; \eqref{hmtfg:escalado:eq:17} proves its uniqueness and the transversality of the crossing. At that root, \eqref{hmtfg:escalado:eq:16} and
\((1+L)(0.001396077+0.00004)<0.001666\)
prove \eqref{hmtfg:escalado:eq:3}.

The bound is uniform in the number of returns: it converts increasing depth into an explicit geometric operator error.


\subsection{Real admissibility of all returns}
\label{hmtfg:escalado:7}

The four initial returns are accompanied, on each rational subinterval, by the inequalities
\begin{equation}
0<a=-f(1)<1,\qquad b=f(a)>a,\qquad f(b)<a.
\label{hmtfg:escalado:eq:19}
\end{equation}
The program stores \(a,b,f(b)\) before each application of \eqref{hmtfg:escalado:eq:7}. Unimodality, evenness and the quadratic critical point are preserved under these normalized returns.

For subsequent iterations, a uniform check in the ball of \eqref{hmtfg:escalado:eq:9} suffices. If \(E=V-V_{g_0}\), the norm \eqref{hmtfg:escalado:eq:5} gives, for \(-0.4\le z\le0\),
\[
|E(z)|\le0.004,\qquad |E'(z)|\le0.01.
\]
The second inequality uses \(\sup_{j\ge1}j(0.4)^{j-1}=1\). Rational evaluation of the central coefficients provides
\begin{equation}
\frac65<V(z)<\frac85,\qquad |V'(z)|<\frac12,
\qquad 0.39<a<0.41.
\label{hmtfg:escalado:eq:20}
\end{equation}
Consequently,
\[
f'(x)=-2x\left[V(z)+\frac25x^2V'(z)\right],
\]
and the bracketed expression exceeds one for \(0\le x\le1\). The map is unimodal and preserves \([-1,1]\). Moreover,
\[
b=f(a)>1-(0.41)^2\frac85=\frac{4569}{6250}>a,
\]
\begin{equation}
f(b)<1-\left(\frac{4569}{6250}\right)^2\frac65
=\frac{35028967}{97656250}<0.39<a.
\label{hmtfg:escalado:eq:21}
\end{equation}
The orbits of the stable leaf satisfy \eqref{hmtfg:escalado:eq:19} at every depth. Thus the analytic iterations of the operator correspond to real doubling returns, with their domains and orientations preserved.


\subsection{Scaling of the local superstable cascade}
\label{hmtfg:escalado:8}

The metric conclusion requires two crossings: that of the family with the stable leaf, proved in section~\ref{hmtfg:escalado:6}, and that of the unstable manifold with the superstable target. The second is proved by Eckmann–Wittwer for
\(\Sigma_2=\{\psi:\psi(1)=0\}\), corresponding to the period-two superstable cycle. Normalization \eqref{hmtfg:escalado:eq:6} agrees with that of the even real operator used there.

The general theorem of Collet–Eckmann–Lanford, section 6 of the cited source, proposition 6.1 and theorems 6.2–6.3, applies to a Banach-space map \(C^2\), a fixed point with a unique unstable eigenvalue \(\delta_{\mathrm F}>1\), a curve transverse to its stable leaf and a target transverse to its unstable manifold. The local preimages of the target intersect the curve at points satisfying
\begin{equation}
\delta_{\mathrm F}^j(\mu_j-\lambda_*)\longrightarrow C\ne0.
\label{hmtfg:escalado:eq:22}
\end{equation}
The index \(j\) counts local returns. Since the initial curve here is \(R^4f_\lambda\), the physical period is \(2^{j+5}\) when the target is taken in \(\Sigma_2\); any fixed number of steps used to bring the target into the local neighborhood is incorporated into a fixed index shift and into \(C\).

The real conditions of section~\ref{hmtfg:escalado:7} and the target’s real doubling branch fix the exact period and itinerary. The sequence is selected by this oriented local continuation. Taking differences in \eqref{hmtfg:escalado:eq:22} gives the ratio limit of \eqref{hmtfg:escalado:eq:4}.

Sources of the two results used: \href{https://link.springer.com/article/10.1007/BF01013368}{Eckmann–Wittwer, \emph{A complete proof of the Feigenbaum conjectures}}; \href{https://people.math.harvard.edu/~knill/history/lanford/papers/ColletEckmannLanford.pdf}{Collet–Eckmann–Lanford, section 6}.

\subsubsection{Power-law asymptotic remainder}

The analyticity of our family allows \eqref{hmtfg:escalado:eq:22} to be made more precise. In the normal-coordinate construction of Collet–Eckmann–Lanford, only the stable manifold, the unstable manifold and the target hypersurface are flattened. These three submanifolds admit local \(C^2\) changes of coordinates with bounded derivatives. The smooth cutoff used is applied in the scalar unstable coordinate; the complete foliation is then obtained through the constructed change of coordinates.

In those coordinates, the correction is written as a series of increments \(w_j\) whose norm \(C^1\) satisfies
\[
\|w_j\|_{C^1}\le C_1\rho^j,\qquad 0<\rho<1.
\]
The chain rule for the cut-off map, with uniformly bounded first- and second-order derivatives, provides
\[
\|w_j\|_{C^2}\le C_2B^j
\]
for some \(B\ge1\). The weighted norm of the construction controls the ordinary \(C^1\) norm on the domain considered. By interpolation,
\[
[Dw_j]_\beta\le C_3
\left(\rho^{1-\beta}B^\beta\right)^j.
\]
Choose
\begin{equation}
0<\beta<\min\left\{1,
\frac{-\log\rho}{\log B-\log\rho}\right\}.
\label{hmtfg:escalado:eq:23}
\end{equation}
The series converges in \(C^{1,\beta}\). Let \(z\) be the resulting normal coordinate, normalized so that the preimages of the target satisfy \(z=\delta_{\mathrm F}^{-j}\), absorbing the fixed index shift. Since our curve is analytic and transverse,
\[
z(\Gamma(\lambda))=c(\lambda-\lambda_*)
+O(|\lambda-\lambda_*|^{1+\beta}),\qquad c\ne0.
\]
Local inversion gives
\begin{equation}
\mu_j-\lambda_*=C\delta_{\mathrm F}^{-j}
+O\left(\delta_{\mathrm F}^{-(1+\beta)j}\right).
\label{hmtfg:escalado:eq:24}
\end{equation}
This proves the remainder in \eqref{hmtfg:escalado:eq:4}, with \(\theta=\delta_{\mathrm F}^{-\beta}\). The statement asserts the existence of the constants in the parameter remainder. The explicit operator rate \(0.83996\) of \eqref{hmtfg:escalado:eq:3}, for its part, controls convergence of the renormalized maps.

\subsubsection{Conversion of the remainder into ratio error}

If \(\mu_j-\lambda_*=C\delta_{\mathrm F}^{-j}(1+e_j)\), with \(|e_j|\le M\theta^j\), write
\[
\mu_{j-1}-\mu_j=C(\delta_{\mathrm F}-1)\delta_{\mathrm F}^{-j}(1+r_j),
\quad
r_j=\frac{\delta_{\mathrm F} e_{j-1}-e_j}{\delta_{\mathrm F}-1}.
\]
With \(K=M(\delta_{\mathrm F}+\theta)/(\delta_{\mathrm F}-1)\), we have \(|r_j|\le K\theta^{j-1}\). Therefore,
\begin{equation}
\boxed{
\left|\frac{\mu_{j-2}-\mu_{j-1}}{\mu_{j-1}-\mu_j}-\delta_{\mathrm F}\right|
\le\frac{\delta_{\mathrm F} K(1+\theta)\theta^{j-2}}
{1-K\theta^{j-1}}}
\label{hmtfg:escalado:eq:25}
\end{equation}
when \(K\theta^{j-1}<1\). This formula separates the evaluation precision of each root from the truncation error of the cascade.


\subsection{Certified continuation of the eight initial roots}
\label{hmtfg:escalado:9}

The finite continuation certificate preserves the family \eqref{hmtfg:escalado:eq:1} and the previous eight-root certificate. For \(2\le n\le8\), it studies
\[
S_n(\lambda)=f_\lambda^{2^n}(0)
\]
from the certified box of \(\lambda_{n-1}\) to \(1.874039\). The complete rational covering contains exactly the inherited root \(\lambda_{n-1}\) and the new root \(\lambda_n\). In particular, \(\lambda_n\) is the unique new root in \((\lambda_{n-1},1.874039]\).

Boxes are excluded by the sign of the return or by a local derivative of certified sign. Anchor boxes retain existence and uniqueness through endpoint signs and local monotonicity. Their endpoints cover each interval exactly; adjacencies are checked. Evaluation uses the moments \eqref{hmtfg:escalado:eq:8}, order-32 Horner evaluation, outer geometric tails and \(|u_0''|\le2\). The uniform tail of the potential is less than \(1.294\,10^{-58}\) on \(|x|\le3/2\).

The final certificate contains 650 processed boxes and 332 covering leaves. The first level has the exact expression \(\lambda_1=1/u_0(1)\). The following approximations are shown solely for guidance; their complete boxes are retained in the certificates:

\begingroup\small
\begin{longtable}{p{0.465\linewidth}p{0.465\linewidth}}
\hline
Level & Superstable parameter \\
\hline
1 & 1.345913990471832792210949095… \\
2 & 1.763379787846910127290794030… \\
3 & 1.850413796950136464530566778… \\
4 & 1.868979756606798125150574309… \\
5 & 1.872955034435659740004205128… \\
6 & 1.873806367467030119412262311… \\
7 & 1.873988695221365362307168780… \\
8 & 1.874027744154450672897007573… \\
\hline\end{longtable}
\endgroup
